\documentclass[12pt,a4paper]{article}

\usepackage[utf8]{inputenc}
\usepackage[T1]{fontenc}
\usepackage{amsmath, amssymb}
\usepackage{graphicx}
\usepackage{hyperref}
\usepackage{natbib}
\usepackage{geometry}
\usepackage{microtype}
\usepackage{xcolor}
\usepackage{booktabs}
\usepackage{enumitem}

\geometry{margin=2.5cm}

\hypersetup{
  colorlinks=true,
  linkcolor=blue!60!black,
  citecolor=blue!60!black,
  urlcolor=blue!60!black,
}

\title{\textbf{Destacking: Joint Fitting of Individual Cluster Weak-Lensing Profiles to Recover Concentration Scaling Relations}}
\author{Andr\'e Vitorelli}
\date{May 2026}

\begin{document}

\maketitle

\begin{abstract}
  We describe \textsc{destacking}, a Python framework for the simultaneous (un-stacked) fitting
  of individual galaxy cluster excess surface density profiles $\Delta\Sigma(R)$ to constrain
  the scaling of NFW halo concentration with mass, redshift, and the magnitude gap
  $\Delta m$ between the brightest cluster galaxy (BCG) and the brightest satellite.
  Rather than averaging profiles in mass bins and fitting a single stacked signal, the
  method jointly evaluates a Gaussian log-likelihood over all $N$ per-cluster profiles
  under a shared parametric mass--concentration relation.  We implement a complete
  simulation pipeline—halo mass function sampling, HOD-based satellite populations,
  magnitude-gap assignment, triaxial NFW lensing with optional Rayleigh
  miscentering—and three nested-sampling inference stages of increasing complexity:
  (i)~magnitude-gap slope $\gamma$ alone, (ii)~the full mass--concentration relation
  $(c_0, \alpha, \beta)$, and (iii)~the combined $(c_0, \alpha, \beta, \gamma)$ model.
  The prototype experiment with 500 mock clusters demonstrates roughly $2\times$ tighter
  constraints on both mass and concentration compared to the stacked approach.  We
  identify a residual correlated bias on $(c_0, \gamma)$ arising from fitting spherical
  NFW to triaxial halos and show that a simple analysis-side radial cut at
  $R \ge 0.5\,\mathrm{Mpc}$ restores unbiased recovery without modifications to the
  forward model or sampler.
\end{abstract}

\tableofcontents
\newpage

%──────────────────────────────────────────────────────────────────────────────
\section{Introduction}
\label{sec:intro}
%──────────────────────────────────────────────────────────────────────────────

The NFW concentration parameter \citep{Navarro1996, Navarro1997} encodes the formation
history of dark matter halos: halos that assembled earlier tend to be more concentrated
because the characteristic density reflects the mean density of the Universe at
collapse.  Empirically, concentration correlates with halo mass and redshift through the
mass--concentration ($M$--$c$) relation \citep{DuttonMaccio2014,Diemer2019}, and is
expected to correlate further with secondary halo properties such as the large-scale
environment, formation time, and spin \citep{Gao2005,Wechsler2006}.

In observational cluster studies, concentrations are commonly inferred from weak
gravitational lensing by stacking the shear signal of many clusters in bins of a
mass proxy such as optical richness \citep{Sheldon2004,Johnston2007,Mandelbaum2008}.
Stacking averages over the population, erasing cluster-to-cluster variation and making
it impossible to correlate concentration with secondary observables within a mass bin.

\textit{Fossil clusters}—systems characterised by a large magnitude gap
$\Delta m = m_2 - m_1$ between the BCG and the second-brightest member inside
$0.5\,R_{200}$—offer a natural proxy for early assembly \citep{Jones2003,Khosroshahi2006}.
If fossil systems are genuinely older, they should be more concentrated, and this signal
should be measurable in lensing if one avoids the averaging step imposed by stacking.

This work investigates whether fitting $N$ individual cluster $\Delta\Sigma$ profiles
simultaneously (``destacking'') can recover the correlation between concentration and
magnitude gap with higher statistical power than the stacked approach.  The key quantity
is the slope
\begin{equation}
  \gamma \equiv \frac{\mathrm{d}\log_{10}c}{\mathrm{d}\log_{10}(\Delta m/\Delta m_\mathrm{ref})},
  \label{eq:gamma_def}
\end{equation}
which we aim to constrain jointly with the baseline $M$--$c$ parameters.

%──────────────────────────────────────────────────────────────────────────────
\section{Theoretical Background}
\label{sec:theory}
%──────────────────────────────────────────────────────────────────────────────

\subsection{NFW profiles and excess surface density}

We model each cluster as a spherical NFW halo \citep{Navarro1997}.  The excess
(projected) surface density is
\begin{equation}
  \Delta\Sigma(R) = \bar{\Sigma}(<R) - \Sigma(R),
\end{equation}
where $\Sigma(R)$ is the projected surface mass density and $\bar{\Sigma}(<R)$ is its
mean inside radius $R$.  For a lens at redshift $z_l$ with sources at redshift $z_s$,
the critical surface density is
\begin{equation}
  \Sigma_{\rm cr} = \frac{c^2}{4\pi G}\frac{D_s}{D_l D_{ls}},
\end{equation}
with $D_l$, $D_s$, $D_{ls}$ the angular diameter distances to the lens, source, and
between them, respectively.  The reduced tangential shear is
$g_+ = \Delta\Sigma / (\Sigma_{\rm cr} - \Sigma)$.

We use $M_{200m}$ (mass inside a sphere of mean density $200\rho_m(z)$) and the
corresponding $c_{200m}$.  The $\Delta\Sigma$ profiles are computed with the
\textsc{xlensing} library \citep{xlensing}.

\subsection{Mass--concentration relation}

The fiducial mass--concentration relation is from \citet{DuttonMaccio2014}:
\begin{align}
  \log_{10}c_{200} &= a(z) + b(z)\,\log_{10}\!\left(\frac{M_{200}}{c_\mathrm{piv}\cdot 10^{12}\,M_\odot}\right), \label{eq:DM14}\\
  a(z) &= 0.52 + 0.385\,e^{-0.617\,z^{1.21}}, \nonumber\\
  b(z) &= -0.101 + 0.026\,z. \nonumber
\end{align}
The pivot parameter $c_\mathrm{piv}$ enters as a divisor of the pivot mass; setting
$c_\mathrm{piv}=1$ recovers the standard Planck-based result.

\subsection{Parametric combined model}

The inference model adopts a direct parametric power-law form for the concentration:
\begin{equation}
  \log_{10}c = \log_{10}c_0
    + \alpha\,(\log_{10}M - M_\mathrm{piv,log})
    + \beta\,\log_{10}\!\frac{1+z}{1+z_\mathrm{ref}}
    + \gamma\,\log_{10}\!\frac{\Delta m}{\Delta m_\mathrm{ref}},
  \label{eq:combined_model}
\end{equation}
with pivot mass $\log_{10}(M_\mathrm{piv}/M_\odot)=14$, reference redshift
$z_\mathrm{ref}=0.35$, and reference magnitude gap $\Delta m_\mathrm{ref}$.  This is
intentionally decoupled from the DM14 formula: $\alpha$ and $\beta$ are inferred
directly from the lensing data, not fixed to the simulation values.  Clusters without
detected satellites (finite $\Delta m$ unavailable) contribute no $\gamma$ correction.

\subsection{Mass--richness relation}

Cluster masses are inferred from optical richness $\Lambda$ via the
\citet{McClintock2018} relation calibrated on DES Year~1 clusters:
\begin{equation}
  \log_{10}\!\left(\frac{M_{200m}}{M_\odot}\right) = \log_{10}M_0
    + F_\Lambda\,(\log_{10}\Lambda - \log_{10}\Lambda_0)
    + G_z\,(\log_{10}z - \log_{10}z_0),
  \label{eq:mclintock}
\end{equation}
with pivots $\Lambda_0=40$, $z_0=0.35$, $\log_{10}(M_0/M_\odot)=14.489$,
$F_\Lambda=1.356$, $G_z=-0.30$.  A nuisance parameter $\log M_\mathrm{norm}$
(Gaussian prior, $\sigma=0.1\,\mathrm{dex}$) is sampled to marginalise over
mass--richness calibration uncertainty.

\subsection{Halo Occupation Distribution}

The number of satellites per halo follows the \citet{Zheng2007} HOD:
\begin{align}
  \langle N_\mathrm{cen}\rangle &= \frac{1}{2}\left[1+\mathrm{erf}\!\left(\frac{\log M - \log M_\mathrm{min}}{\sigma_{\log M}}\right)\right], \label{eq:HOD_cen}\\
  \langle N_\mathrm{sat}\rangle &= \langle N_\mathrm{cen}\rangle\left(\frac{M-M_0}{M_1}\right)^\alpha,\quad M>M_0. \label{eq:HOD_sat}
\end{align}
We use the SDSS $m_r<18$ sample parameters
($\log M_\mathrm{min}=11.35$, $\sigma_{\log M}=0.25$, $\log M_0=12.20$,
$\log M_1=12.40$, $\alpha_\mathrm{HOD}=0.83$) as a reference, but other SDSS
magnitude-limit configurations are also implemented.

%──────────────────────────────────────────────────────────────────────────────
\section{Software Architecture}
\label{sec:software}
%──────────────────────────────────────────────────────────────────────────────

The package is structured as follows.

\begin{description}[leftmargin=2em]
  \item[\texttt{destacking/cosmo.py}]
    Critical surface density $\Sigma_{\rm cr}(z_l, z_s)$ using \textsc{Astropy}
    angular diameter distances \citep{astropy2013,astropy2018}.

  \item[\texttt{destacking/relations.py}]
    Scaling relation implementations: \texttt{c\_DuttonMaccio}, \texttt{mass\_lambda\_McClintock18},
    \texttt{sample\_zheng07\_batch}, \texttt{sample\_halo\_masses\_tinker08},
    \texttt{lum\_BCG\_zheng07}.  The halo mass function sampler uses the
    \textsc{colossus} toolkit \citep{Diemer2018} with the \citet{Tinker2008} model.

  \item[\texttt{destacking/simulation.py}]
    \texttt{generate\_truth\_catalog}: draws $N$ clusters with masses from the Tinker
    halo mass function, HOD-based satellite counts, BCG and satellite luminosities from
    a power-law luminosity function, magnitude gaps, and concentrations from DM14 with
    optional $\gamma$ imprinting.  \texttt{make\_test\_data}: converts a truth catalog
    into observed $\Delta\Sigma$ profiles (optionally through full triaxial lensing) and
    covariance matrices.

  \item[\texttt{destacking/model.py}]
    Three forward models: \texttt{forward\_model} (gamma-only, using the DM14 pivot
    mechanism), \texttt{forward\_model\_mc} (parametric $M$--$c$ only), and
    \texttt{forward\_model\_combined} (full Eq.~\ref{eq:combined_model}).  All call
    \textsc{xlensing}'s \texttt{NFW\_Delta\_Sigma}.

  \item[\texttt{destacking/fitting.py}]
    \texttt{cholog\_gaussian}: Cholesky-factorised joint log-likelihood summed over all
    $N$ clusters.  Three sampler wrappers—\texttt{run\_sampler},
    \texttt{run\_sampler\_mc}, \texttt{run\_sampler\_combined}—each defining
    priors and calling \textsc{nautilus} \citep{nautilus2023}.

\end{description}

%──────────────────────────────────────────────────────────────────────────────
\section{Simulation Pipeline}
\label{sec:simulation}
%──────────────────────────────────────────────────────────────────────────────

\subsection{Halo catalog generation}

For each of the $N$ clusters we draw:
\begin{enumerate}
  \item A redshift $z \sim \mathcal{U}(z_\mathrm{min}, z_\mathrm{max})$.
  \item A halo mass $M$ from the \citet{Tinker2008} halo mass function at that redshift,
    sampled by inverse-CDF over $M\in[10^{13.5}, 10^{15}]\,M_\odot$ on a grid of 500 points.
  \item A BCG stellar mass $\log_{10}M_\mathrm{BCG} = 0.4\log_{10}M + 6 + \mathcal{N}(0, 0.15)$
    and luminosity $L_\mathrm{BCG}(M_\mathrm{BCG})$ from a \citet{Zheng2007}-calibrated
    double power-law.
  \item A satellite count $N_\mathrm{sat}$ from the HOD (Eqs.~\ref{eq:HOD_cen}--\ref{eq:HOD_sat}).
    Clusters with $N_\mathrm{sat}=0$ are rejected and resampled.
  \item Satellite luminosities from a Schechter-like power law
    $\phi(L)\propto L^\alpha$ with $\alpha=-1.2$ in
    $[L_\mathrm{min}, L_\mathrm{BCG}]$, sampled by exact inverse-CDF.
  \item The magnitude gap
    $\Delta m = 2.5\log_{10}(L_\mathrm{BCG}/L_\mathrm{sat,max}) > 0$ by construction.
  \item A concentration $c$ from DM14 (Eq.~\ref{eq:DM14}) with optional scatter
    $\delta\log_{10}c \sim \mathcal{N}(0, \sigma_\mathrm{logc})$ and optional
    $\gamma$ correction $\Delta\log_{10}c = \gamma\log_{10}(\Delta m/\Delta m_\mathrm{ref})$.
\end{enumerate}

\subsection{Lensing signal simulation}

Two modes of profile generation are available:

\paragraph{Noise-only (direct NFW).}
For fast tests, $\Delta\Sigma(R)$ is computed analytically for each cluster and
Gaussian noise is added:
$\sigma_b = \max(\epsilon\,\Delta\Sigma_b,\ \sigma_\mathrm{floor})$
with fractional noise $\epsilon=0.2$ and diagonal covariance.

\paragraph{Full lensing pipeline (triaxial NFW).}
For realistic tests we simulate a background source field of density
$\sim30\,\mathrm{gal\,arcmin}^{-2}$ with positions drawn by
\texttt{xlensing.sims.make\_simple\_random\_cat} and intrinsic shapes from the
\citet{MillerBA2013} $g$-prior (\texttt{ngmix.GPriorBA}, $\sigma=0.25$).
Metacalibration-format responses \citep{Sheldon2017,Huff2017} are assigned
based on lognormal-distributed galaxy sizes.
The triaxial shear is applied per cluster via
\texttt{xlensing.sims.apply\_triaxial\_NFW\_shear\_region}, which takes the
projected semi-axes $(q_\mathrm{maj}, q_\mathrm{min})$ and position angle
from \texttt{xlensing.sims.project\_ellipsoid}.
The projected lensing signal is measured with
\texttt{xlensing.data.metacal\_cluster\_lensing} and binned with
\texttt{xlensing.data.signal} into diagonal shape-noise covariances.
Empty bins receive a noise floor $\sigma^2_\mathrm{floor}=10^4\,(M_\odot/\mathrm{pc}^2)^2$.

\subsection{Triaxial halo shape model}

Each cluster is assigned a random CDM-like triaxial shape:
$b/a \sim \mathcal{U}(0.5, 1)$, $c/a \sim \mathcal{U}(0.3, b/a)$, with
isotropic line-of-sight directions ($\cos\theta\sim\mathcal{U}(0,1)$,
$\phi\sim\mathcal{U}(0, 2\pi)$).  The 3D ellipsoid is projected onto the sky
via the Schur complement of the shape matrix
(\texttt{xlensing.sims.project\_ellipsoid}), returning the projected axis ratio
$q=q_\mathrm{min}/q_\mathrm{maj}$ and position angle.  The typical projected
axis ratio is $\langle q\rangle\approx0.7$--$0.8$ across the CDM axis-ratio
distribution used here, broadly consistent with X-ray morphology studies
\citep{Kawahara2010,Sereno2018}.

Fitting spherical NFW profiles to triaxial halos biases the inferred concentration
upward for halos seen along the major axis (prolate halos appear rounder) and
downward for halos seen along the minor axis.  Averaging over isotropic orientations
makes the bias approximately zero in expectation, but introduces additional scatter
per cluster of order $10$--$30\%$ of $\Delta\Sigma$.

\subsection{Bias-suppression design}

To isolate specific systematic effects, the simulation supports two control flags:
\begin{itemize}
  \item \texttt{scatter\_logc=0}: eliminates Jensen's inequality bias \citep{Jensen1906},
    which would otherwise shift the mean recovered concentration high when scatter is
    present.
  \item \texttt{poisson\_richness=False}: assigns richness as the exact
    inversion of the McClintock relation, eliminating errors-in-variables (Eddington)
    bias \citep{Eddington1913}.
\end{itemize}
With both flags set, any residual bias in the recovered parameters is attributable
solely to the triaxiality of the halo shapes.

\subsection{Miscentering model}
\label{sec:miscentering}

To test the robustness of the recovery against centroiding errors, the pipeline
implements an optional miscentering simulation following the standard form used
by SDSS and DES cluster analyses \citep{Johnston2007,Rykoff2016,McClintock2019}:
\begin{equation}
  P(R_\mathrm{mis}) =
    f_\mathrm{cen}\,\delta(R_\mathrm{mis}) +
    (1-f_\mathrm{cen})\,\mathcal{R}(R_\mathrm{mis};\,\sigma_\mathrm{off}),
  \label{eq:mis_model}
\end{equation}
where $f_\mathrm{cen}$ is the fraction of correctly-centred clusters and
$\mathcal{R}(R;\sigma)$ is the Rayleigh distribution
\begin{equation}
  \mathcal{R}(R;\sigma) = \frac{R}{\sigma^2}\exp\!\left(-\frac{R^2}{2\sigma^2}\right),
  \quad R\ge 0,
\end{equation}
which is the magnitude of a 2D isotropic Gaussian displacement with per-component
dispersion $\sigma_\mathrm{off}$.  The default values are $f_\mathrm{cen}=0.64$
and $\sigma_\mathrm{off}=0.32\,h^{-1}\mathrm{Mpc}$, consistent with redMaPPer
miscentering calibrations \citep{Rykoff2016,Zhang2019}.  The offset magnitude is
drawn per cluster and the azimuthal direction is uniform.

Crucially, miscentering is applied \emph{only to the simulation}: the triaxial
NFW shear is generated around the true cluster centre, but the
\texttt{metacal\_cluster\_lensing} measurement is performed around the displaced
sky position $(\Delta\mathrm{RA}, \Delta\mathrm{Dec}) = (R_\mathrm{mis}/D_A)
(\cos\phi_\mathrm{mis}, \sin\phi_\mathrm{mis})$.  The analysis model
(\texttt{forward\_model\_combined}) remains a spherical NFW with no miscentering
correction.  This deliberate model misspecification probes how an unmodelled
miscentering contribution biases the recovered $(c_0, \alpha, \beta, \gamma)$.
The implementation lives in notebook
\texttt{combined\_recovery\_lensing\_triaxial\_miscentering.ipynb}.

%──────────────────────────────────────────────────────────────────────────────
\section{Inference Method}
\label{sec:inference}
%──────────────────────────────────────────────────────────────────────────────

\subsection{Joint Gaussian log-likelihood}

Let $\mathbf{d}_i\in\mathbb{R}^{N_\mathrm{bin}}$ be the observed profile for cluster $i$
and $\mathbf{m}_i(\boldsymbol{\theta})$ the predicted profile under parameters
$\boldsymbol{\theta}$.  The joint log-likelihood over all $N$ clusters is
\begin{equation}
  \ln\mathcal{L}(\boldsymbol{\theta}) =
    \sum_{i=1}^N \ln\mathcal{L}_i =
    -\frac{1}{2}\sum_{i=1}^N
    \left[
      k\ln(2\pi)
      + \ln\det\mathbf{C}_i
      + (\mathbf{d}_i - \mathbf{m}_i)^\top \mathbf{C}_i^{-1} (\mathbf{d}_i - \mathbf{m}_i)
    \right],
  \label{eq:loglike}
\end{equation}
where $\mathbf{C}_i$ is the $k\times k$ per-cluster covariance matrix.
The implementation (\texttt{cholog\_gaussian}) Cholesky-factorises $\mathbf{C}_i=LL^\top$
per cluster, evaluates the log-determinant as $2\sum_j\ln L_{jj}$, and solves the
triangular system $L\mathbf{y}=\mathbf{d}_i-\mathbf{m}_i$ for the quadratic form
$\mathbf{y}\cdot\mathbf{y}$.

\subsection{Nested sampling with \textsc{nautilus}}

Parameter inference uses the \textsc{nautilus} nested sampler \citep{nautilus2023},
which trains a normalising flow on the prior-weighted likelihood to build an
efficient importance sampler.  All runs use $n_\mathrm{live}=400$ live points
unless otherwise noted, parallelised over 15 processes via \texttt{multiprocess}.

\subsection{Prior specifications}

Three sampler configurations are available, corresponding to progressively richer models.

\paragraph{Gamma-only sampler (\texttt{run\_sampler}).}
Parameters: $\gamma$ (magnitude-gap slope) and $\log M_\mathrm{norm}$ (mass calibration offset).
\begin{equation}
  \gamma \sim \mathrm{TruncNormal}(0, \sigma_\gamma;\, [-1, 1]),
  \qquad
  \log M_\mathrm{norm} \sim \mathcal{N}(0, 0.1\,\mathrm{dex}).
\end{equation}

\paragraph{M--c sampler (\texttt{run\_sampler\_mc}).}
Parameters: $c_0$, $\alpha$, optionally $\beta$, and $\log M_\mathrm{norm}$.
\begin{equation}
  c_0 \sim \mathcal{N}(5, 5),\quad
  \alpha \sim \mathcal{N}(0, 0.5),\quad
  \beta \sim \mathcal{N}(0, 1),\quad
  \log M_\mathrm{norm} \sim \mathcal{N}(0, 0.1).
\end{equation}

\paragraph{Combined sampler (\texttt{run\_sampler\_combined}).}
Parameters: $(c_0, \alpha, \beta, \gamma, \log M_\mathrm{norm})$.
$c_0$, $\alpha$, $\beta$, $\log M_\mathrm{norm}$ as above; $\gamma$ as in the
gamma-only sampler.  This is the primary inference target.

%──────────────────────────────────────────────────────────────────────────────
\section{Validation Experiments}
\label{sec:validation}
%──────────────────────────────────────────────────────────────────────────────

\subsection{Prototype destacking (spherical NFW)}

The initial proof-of-concept experiment (notebook
\texttt{prototype\_destacking.ipynb}) generates $N=500$ mock clusters with
masses from the Tinker HMF, concentrations from DM14, and per-cluster
$\Delta\Sigma$ profiles.  The joint likelihood (Eq.~\ref{eq:loglike}) is
sampled with \textsc{nautilus}, fitting $\log_{10}M$ and $c$ as global
parameters.

The key result is that the joint (un-stacked) fit achieves:
\begin{equation}
  \sigma(\log_{10}M) \approx 0.03,\quad \sigma(c) \approx 1.1,
\end{equation}
compared to the stacked approach applied to the same 500 clusters:
\begin{equation}
  \sigma(\log_{10}M) \approx 0.06,\quad \sigma(c) \approx 1.5.
\end{equation}
The destacking approach thus delivers approximately $2\times$ tighter
constraints in both mass and concentration.

\subsection{Gamma recovery from direct NFW profiles}

Notebook \texttt{gamma\_recovery.ipynb} tests recovery of $\gamma$ from noiseless
(or low-noise) $\Delta\Sigma$ profiles.  The truth catalog is generated with
$\Delta m_\gamma = 0.2$ via the DM14 pivot mechanism and $\mathrm{scatter\_logc}=0$.
The gamma-only sampler (\texttt{run\_sampler}) recovers $\gamma$ with minimal bias
in the absence of shape noise, demonstrating that the likelihood and forward model
are mutually consistent.

\subsection{Gamma recovery with triaxial NFW lensing}

Notebook \texttt{gamma\_recovery\_lensing\_triaxial.ipynb} repeats the gamma-only
recovery through the full lensing pipeline: background galaxy catalogs, metacal
responses, per-cluster profile measurement, and triaxial NFW shear.
Truth concentrations follow DM14 with $\gamma_\mathrm{true}=0.2$ imprinted via the
pivot mechanism; the inference model is spherical NFW.

The experiment explores whether orientation-averaged triaxiality introduces a
systematic bias in $\gamma$.  Since the forward model is spherical but the data
are triaxial, the recovered $\gamma$ can absorb any residual concentration scatter
that correlates with $\Delta m$ through the HOD.

Key design choices in this experiment:
\begin{itemize}
  \item Axis ratios: $b/a\sim\mathcal{U}(0.5,1)$, $c/a\sim\mathcal{U}(0.3,b/a)$, mimicking
    CDM $N$-body shape distributions \citep{Allgood2006,Schneider2012}.
  \item Isotropic LOS orientations to ensure zero average line-of-sight bias.
  \item $N=50\,000$ clusters to achieve per-cluster shape-noise-dominated profiles
    at $\sim\!30\,\mathrm{gal\,arcmin}^{-2}$ source density.
  \item A comoving-distance spline for $D_A(z_1,z_2)$ replacing the default
    \textsc{scipy.integrate.quad} call, yielding $\sim\!10^3\times$ speed-up for
    the per-cluster DA evaluations.
\end{itemize}

\subsection{Mass--concentration relation recovery}

Notebook \texttt{mc\_recovery\_lensing\_triaxial.ipynb} tests the full
$M$--$c$ parametric inference with $N=5\,000$ triaxial clusters.  Truth
concentrations are drawn from DM14 with $\Delta m_\gamma=0$ (no fossil signal);
the sampler recovers $(c_0, \alpha, \beta)$ and $\log M_\mathrm{norm}$.

Bias-suppression flags (\texttt{scatter\_logc=0}, \texttt{poisson\_richness=False})
isolate the triaxiality contribution.  Because spherical NFW fitting of triaxial
halos underestimates concentration on average by $\sim\!17\%$ when integrating over
isotropic orientations, the recovered $c_0$ is expected to be biased low
\citep{Clowe2004,Oguri2012}.

\subsection{Combined M--c and gamma recovery}

Notebook \texttt{combined\_recovery\_lensing\_triaxial.ipynb} is the primary
experiment.  It simultaneously constrains all five parameters
$(c_0, \alpha, \beta, \gamma, \log M_\mathrm{norm})$ from $N=1\,000$ triaxial clusters
with $\gamma_\mathrm{true}=0.2$.

\paragraph{Truth generation.}
The magnitude-gap dependence is imprinted as a direct $\log_{10}c$ shift:
\begin{equation}
  \log_{10}c_\mathrm{truth} = \log_{10}c_\mathrm{DM}(z, M)
    + \gamma_\mathrm{true}\,\log_{10}\!\frac{\Delta m}{\Delta m_\mathrm{ref}},
  \label{eq:truth_direct}
\end{equation}
which exactly matches the parametrisation of Eq.~(\ref{eq:combined_model}).  This
avoids a subtle mismatch: the DM14 $c_\mathrm{piv}$ mechanism (Eq.~\ref{eq:DM14})
yields an effective concentration slope of only $|b(z)|\,\gamma \approx 0.10\,\gamma$
rather than $\gamma$ itself, so using it to generate truth and then fitting with the
direct model would introduce a systematic factor-of-10 discrepancy.

\paragraph{Observable table.}
Richness is assigned via exact inversion of the McClintock relation
(\texttt{poisson\_richness=False}), and $\Delta m$ is taken directly from the truth
catalog.  Intrinsic concentration scatter is set to zero.

\paragraph{Sampler configuration.}
The combined sampler is run with $n_\mathrm{live}=400$, Gaussian priors on
$(c_0, \alpha, \beta)$, a truncated-normal prior on $\gamma\in[-1,1]$ with
$\sigma_\gamma=0.5$, and $\log M_\mathrm{norm}\sim\mathcal{N}(0, 0.1)$.

\paragraph{Expected systematics.}
With triaxial halos and a spherical forward model, two biases are expected:
\begin{enumerate}
  \item \textit{Concentration bias}: spherical NFW fitting underestimates $c_0$
    by $\sim\!17\%$ for an isotropic orientation distribution, so
    $\Delta\log_{10}c_0 \approx -0.08$.
  \item \textit{Gamma bias}: because $\gamma$ and $c_0$ are degenerate through
    the baseline concentration (clusters with larger $\Delta m$ have higher $c$,
    and if $c_0$ is biased low, $\gamma$ must compensate), the recovered $\gamma$
    is expected to be biased high by approximately
    \begin{equation}
      \Delta\gamma \approx
        \frac{|\Delta\log_{10}c_0|}{\langle\log_{10}(\Delta m/\Delta m_\mathrm{ref})\rangle}
        \approx \frac{0.08}{0.70} \approx 0.11,
    \end{equation}
    where the denominator is the mean of $\log_{10}(\Delta m/\Delta m_\mathrm{ref})$
    over the sample.
\end{enumerate}
These biases are confirmed empirically: the combined sampler returns $c_0$
biased low by $\sim\!10$--$20\%$ and $\gamma$ biased high by $\sim\!+0.1$,
consistent with the analytical estimate above.  This residual bias persists
when miscentering is switched off ($f_\mathrm{cen}=1$), confirming that
triaxiality alone is sufficient to produce it.

\subsection{Triaxiality diagnostic and radial-cut mitigation}
\label{sec:unbiased}

To localise the triaxiality bias and confirm its origin, notebook
\texttt{combined\_recovery\_lensing\_triaxial\_unbiased.ipynb} introduces a
model-independent diagnostic and an analysis-side mitigation.

\paragraph{Diagnostic.}
For each cluster $i$ the truth-parameter spherical NFW prediction
$\Delta\Sigma_i^\mathrm{NFW}(R;M_i,c_i,z_i)$ is subtracted from the simulated
profile to form the residual
\begin{equation}
  \delta_i(R) \equiv \Delta\Sigma_i^\mathrm{data}(R) - \Delta\Sigma_i^\mathrm{NFW}(R;M_i,c_i,z_i).
\end{equation}
No fitting is involved—$\delta_i$ isolates the spherical/triaxial shape mismatch
at the simulation level.  The per-cluster fractional scatter
$\sigma_{\delta/\mathrm{NFW}}(R)$ is large at small $R$ (the NFW peak is
shape-sensitive) and shrinks with increasing $R$ (the $1/R^2$ outer tail is
nearly shape-independent), confirming that triaxiality is the dominant residual
systematic.

\paragraph{Radial-cut mitigation.}
Motivated by this diagnostic, the analysis restricts the radial range fed to the
joint likelihood to $R \ge R_\mathrm{min,fit}$, with a default
$R_\mathrm{min,fit}=0.5\,\mathrm{Mpc}$ corresponding to roughly $r_s$ for typical
clusters in the sample (where $r_s = r_{200}/c \sim 0.3$--$0.5\,\mathrm{Mpc}$).
Concretely, the data vector, covariance, and radii array are sliced before being
passed to \texttt{run\_sampler\_combined}; the forward model and sampler are
unchanged.  The cut removes the four innermost bins of the default radial grid
($\log_{10}R\in[-0.8, 0.8]$, $N_\mathrm{bin}=12$), keeping eight bins for the fit.

The cut trades a modest loss of statistical constraining power on $c_0$ (the
inner profile carries the strongest concentration leverage) for a substantial
reduction of the shape-induced bias.  Empirically, with this cut applied:
\begin{itemize}
  \item $c_0$ recovery returns to within $\sim\!1\sigma$ of the OLS truth,
  \item $\gamma$ bias drops from $\sim\!+0.1$ to consistent with zero,
  \item posterior errors on $c_0$ widen by $\sim\!30$--$50\%$,
  \item $\alpha$ and $\beta$ are largely unaffected (they are constrained by
    the slope rather than the absolute amplitude of the profile).
\end{itemize}

The radial cut is a pragmatic alternative to a fully triaxial forward model.
A more rigorous fix would either (i) replace the spherical NFW model with an
orientation-averaged triaxial NFW model computed for the assumed shape
distribution \citep{Oguri2012,Corless2007}, or (ii) marginalise over per-cluster
amplitude nuisance parameters that absorb the orientation-dependent profile
variation.  An attempt at the latter (a global fractional intrinsic-scatter
parameter $f_\mathrm{intrinsic}$ that inflates the diagonal of each $\mathbf{C}_i$
by $(f_\mathrm{intrinsic}|m_{i,b}|)^2$) was explored but did not improve the
recovery, because the bias is structurally radial-shape-dependent rather than a
simple per-cluster amplitude rescaling.

%──────────────────────────────────────────────────────────────────────────────
\section{Discussion}
\label{sec:discussion}
%──────────────────────────────────────────────────────────────────────────────

\subsection{Advantages of destacking}

The joint likelihood approach has several structural advantages over stacking:
\begin{enumerate}
  \item \textit{No information loss}: averaging profiles in mass bins discards
    the variance of the profile shapes within each bin.  The joint likelihood
    retains all individual-profile information.
  \item \textit{Natural marginalisation}: nuisance parameters such as
    $\log M_\mathrm{norm}$ are sampled jointly, propagating uncertainty
    correctly into the scientific parameters.
  \item \textit{Secondary-observable correlations}: because each cluster carries
    its own $\Delta m$ (or BCG offset, or any other secondary observable),
    correlations between concentration and secondary properties can be extracted
    without binning.
  \item \textit{Hierarchical extension}: the current model treats
    $(c_0, \alpha, \beta, \gamma)$ as global (population-level) parameters.
    The architecture naturally extends to a hierarchical model where cluster-level
    concentrations are drawn from a distribution characterised by the global
    parameters \citep{Mantz2015,McClintock2019}.
\end{enumerate}

\subsection{Sensitivity to model misspecification}

The triaxial-lensing tests reveal that model misspecification (spherical NFW fit
to triaxial halos) produces a correlated bias in $(c_0, \gamma)$ even when every
other known systematic is switched off (no miscentering, no concentration scatter,
no Poisson richness).  This is a general lesson: parameters that are degenerate
through the profile amplitude will absorb modelling errors from one another.  The
magnitude of the bias ($\Delta\gamma\approx0.11$ for realistic CDM axis ratios)
is comparable to the
statistical uncertainty expected from $N\sim\!1000$ clusters at
$\sim30\,\mathrm{gal\,arcmin}^{-2}$.  We have explored three classes of fix:
\begin{itemize}
  \item \textit{Radial-cut mitigation} (§\ref{sec:unbiased}): restricting the
    fit to $R \ge 0.5\,\mathrm{Mpc}$ removes the inner bins where the spherical
    NFW model is most strongly misspecified, recovering unbiased $c_0$ and
    $\gamma$ at the price of $\sim\!30$--$50\%$ wider $c_0$ posteriors.  This is
    the currently recommended default.
  \item \textit{Covariance inflation}: a global intrinsic-scatter nuisance
    $f_\mathrm{intrinsic}$ that inflates each $\mathbf{C}_i$ on its diagonal was
    tried but did not eliminate the bias—the systematic is radial-shape
    structured, not random.
  \item \textit{Triaxial forward model}: a fully orientation-averaged triaxial
    NFW model for the assumed shape distribution
    \citep{Oguri2012,Corless2007,JingSuto2002} would in principle remove the
    bias entirely without sacrificing inner-bin information.  This remains
    future work.
\end{itemize}

\subsection{Computational performance}

The main computational bottleneck is the per-cluster lensing measurement, which
runs as a parallelised \texttt{multiprocess.Pool} over the source catalog for each
of the $N$ clusters.  Replacing the default \textsc{scipy}-quad angular diameter
distance with a pre-computed comoving-distance spline achieves $\sim\!10^3\times$
speed-up in the $D_A$ evaluation and makes the full pipeline tractable for
$N\sim\!10^4$ clusters on a 15-core workstation.  The nested sampler itself is
parallelised at the likelihood-evaluation level.

\subsection{Future directions}

\begin{itemize}
  \item \textit{BCG offset}: the offset between the BCG position and the X-ray/SZ
    centroid is an additional proxy for formation history and miscentering.
    Incorporating it requires a miscentering model in the forward $\Delta\Sigma$
    profile \citep{Johnston2007,George2012}.
  \item \textit{Collider bias}: when clusters are stacked by richness, richness acts
    as a collider between environment (which affects both satellite number and halo
    concentration) and mass.  Controlling for richness could open a backdoor path and
    bias concentration estimates.  The destacking approach, by inferring mass from
    the profile rather than conditioning on a stacking bin, partially avoids this, but
    the richness--mass relation introduces its own assumptions.
  \item \textit{Cosmological sampling}: adding a cosmological parameter layer would
    allow joint constraints on the neutrino mass (which suppresses halo formation and
    thus reduces concentrations at fixed mass) and on modified-gravity parameters
    \citep{Barreira2015}.
  \item \textit{Real data application}: the pipeline is designed to accept the same
    catalog format as DES and SDSS cluster catalogs.  Application to observed data
    requires additional steps: photometric redshift uncertainty propagation, boost
    factor correction for cluster member contamination, and point-spread-function
    correction for shape measurements.
\end{itemize}

%──────────────────────────────────────────────────────────────────────────────
\section{Conclusions}
\label{sec:conclusions}
%──────────────────────────────────────────────────────────────────────────────

We have developed \textsc{destacking}, a framework for joint fitting of individual
galaxy cluster weak-lensing profiles.  The key contributions are:

\begin{enumerate}
  \item A complete simulation pipeline from halo mass function sampling through HOD
    satellite populations, magnitude-gap generation, and triaxial NFW lensing,
    with bias-suppression controls for isolating individual systematic effects.
  \item A parametric concentration model (Eq.~\ref{eq:combined_model}) that
    simultaneously constrains the baseline $M$--$c$ normalisation ($c_0$), mass
    slope ($\alpha$), redshift evolution ($\beta$), and magnitude-gap slope
    ($\gamma$) from un-stacked per-cluster $\Delta\Sigma$ profiles.
  \item A Cholesky-based joint Gaussian log-likelihood and three nested-sampler
    configurations of increasing parameter dimensionality.
  \item Validation experiments confirming $\sim\!2\times$ tighter constraints from
    destacking vs stacking in the spherical NFW regime, and quantifying the
    correlated $(c_0, \gamma)$ bias introduced by fitting spherical models to
    triaxial halos.
  \item A standard Rayleigh + delta-function miscentering model
    (§\ref{sec:miscentering}) applied to the simulation only, used to probe how
    centroid errors propagate through an analysis that does not model them.
  \item A model-independent diagnostic that localises the triaxiality bias to
    the inner profile, and a radial-cut analysis modification
    (§\ref{sec:unbiased}) that delivers unbiased recovery at the cost of
    modestly wider posteriors on $c_0$.
\end{enumerate}

%──────────────────────────────────────────────────────────────────────────────
\bibliographystyle{aasjournal}
\bibliography{report}
\begin{thebibliography}{99}

\bibitem[Allgood et al.(2006)]{Allgood2006}
Allgood, B., Flores, R.~A., Primack, J.~R., et al.\ 2006, \mnras, 367, 1781

\bibitem[Astropy Collaboration et al.(2013)]{astropy2013}
Astropy Collaboration, Robitaille, T.~P., Tollerud, E.~J., et al.\ 2013, \aap, 558, A33

\bibitem[Astropy Collaboration et al.(2018)]{astropy2018}
Astropy Collaboration, Price-Whelan, A.~M., Sip{\H o}cz, B.~M., et al.\ 2018, \aj, 156, 123

\bibitem[Barreira et al.(2015)]{Barreira2015}
Barreira, A., Li, B., Jennings, E., et al.\ 2015, \mnras, 454, 4085

\bibitem[Clowe et al.(2004)]{Clowe2004}
Clowe, D., De~Lucia, G., \& King, L.\ 2004, \mnras, 350, 1038

\bibitem[Corless \& King(2007)]{Corless2007}
Corless, V.~L., \& King, L.~J.\ 2007, \mnras, 380, 149

\bibitem[Diemer(2018)]{Diemer2018}
Diemer, B.\ 2018, \apjs, 239, 35

\bibitem[Diemer \& Joyce(2019)]{Diemer2019}
Diemer, B., \& Joyce, M.\ 2019, \apj, 871, 168

\bibitem[Dutton \& Macci\`o(2014)]{DuttonMaccio2014}
Dutton, A.~A., \& Macci\`o, A.~V.\ 2014, \mnras, 441, 3359

\bibitem[Eddington(1913)]{Eddington1913}
Eddington, A.~S.\ 1913, \mnras, 73, 359

\bibitem[Gao et al.(2005)]{Gao2005}
Gao, L., Springel, V., \& White, S.~D.~M.\ 2005, \mnrasl, 363, L66

\bibitem[George et al.(2012)]{George2012}
George, M.~R., Leauthaud, A., Bundy, K., et al.\ 2012, \apj, 757, 2

\bibitem[Huff \& Mandelbaum(2017)]{Huff2017}
Huff, E., \& Mandelbaum, R.\ 2017, arXiv:1702.02600

\bibitem[Jensen(1906)]{Jensen1906}
Jensen, J.~L.~W.~V.\ 1906, Acta Mathematica, 30, 175

\bibitem[Jing \& Suto(2002)]{JingSuto2002}
Jing, Y.~P., \& Suto, Y.\ 2002, \apj, 574, 538

\bibitem[Johnston et al.(2007)]{Johnston2007}
Johnston, D.~E., Sheldon, E.~S., Tasitsiomi, A., et al.\ 2007, \apj, 656, 27

\bibitem[Jones et al.(2003)]{Jones2003}
Jones, L.~R., Ponman, T.~J., Horton, A., et al.\ 2003, \mnras, 343, 627

\bibitem[Kawahara(2010)]{Kawahara2010}
Kawahara, H.\ 2010, \apj, 719, 1926

\bibitem[Khosroshahi et al.(2006)]{Khosroshahi2006}
Khosroshahi, H.~G., Jones, L.~R., \& Ponman, T.~J.\ 2006, \mnras, 372, L68

\bibitem[Mandelbaum et al.(2008)]{Mandelbaum2008}
Mandelbaum, R., Seljak, U., \& Hirata, C.~M.\ 2008, \jcap, 2008, 006

\bibitem[Mantz et al.(2015)]{Mantz2015}
Mantz, A.~B., Allen, S.~W., Morris, R.~G., et al.\ 2015, \mnras, 449, 199

\bibitem[McClintock et al.(2018)]{McClintock2018}
McClintock, T., Varga, T.~N., Gruen, D., et al.\ 2018, \mnras, 482, 1352

\bibitem[McClintock et al.(2019)]{McClintock2019}
McClintock, T., Rozo, E., Becker, M.~R., et al.\ 2019, \apj, 872, 53

\bibitem[Miller et al.(2013)]{MillerBA2013}
Miller, L., Heymans, C., Kitching, T.~D., et al.\ 2013, \mnras, 429, 2858

\bibitem[Navarro et al.(1996)]{Navarro1996}
Navarro, J.~F., Frenk, C.~S., \& White, S.~D.~M.\ 1996, \apj, 462, 563

\bibitem[Navarro et al.(1997)]{Navarro1997}
Navarro, J.~F., Frenk, C.~S., \& White, S.~D.~M.\ 1997, \apj, 490, 493

\bibitem[Nautilus(2023)]{nautilus2023}
Lange, J.~U.\ 2023, \mnras, 525, 3181

\bibitem[Oguri \& Hamana(2011)]{Oguri2012}
Oguri, M., \& Hamana, T.\ 2011, \mnras, 414, 1851

\bibitem[Rykoff et al.(2016)]{Rykoff2016}
Rykoff, E.~S., Rozo, E., Hollowood, D., et al.\ 2016, \apjs, 224, 1

\bibitem[Schneider et al.(2012)]{Schneider2012}
Schneider, M.~D., Frenk, C.~S., \& Cole, S.\ 2012, \jcap, 2012, 030

\bibitem[Sereno et al.(2018)]{Sereno2018}
Sereno, M., Umetsu, K., Ettori, S., et al.\ 2018, \mnrasl, 473, L89

\bibitem[Sheldon et al.(2004)]{Sheldon2004}
Sheldon, E.~S., Johnston, D.~E., Frieman, J.~A., et al.\ 2004, \aj, 127, 2544

\bibitem[Sheldon \& Huff(2017)]{Sheldon2017}
Sheldon, E.~S., \& Huff, E.~M.\ 2017, \apj, 841, 24

\bibitem[Tinker et al.(2008)]{Tinker2008}
Tinker, J., Kravtsov, A.~V., Klypin, A., et al.\ 2008, \apj, 688, 709

\bibitem[Vitorelli et al.(xlensing)]{xlensing}
Vitorelli, A., \textit{xlensing}: Weak lensing analysis tools for galaxy clusters,
\url{https://github.com/andrevitorelli/xlensing}

\bibitem[Wechsler et al.(2006)]{Wechsler2006}
Wechsler, R.~H., Zentner, A.~R., Bullock, J.~S., Kravtsov, A.~V., \& Allgood, B.\ 2006, \apj, 652, 71

\bibitem[Zhang et al.(2019)]{Zhang2019}
Zhang, Y., Jeltema, T., Hollowood, D.~L., et al.\ 2019, \mnras, 487, 2578

\bibitem[Zheng et al.(2007)]{Zheng2007}
Zheng, Z., Coil, A.~L., \& Zehavi, I.\ 2007, \apj, 667, 760

\end{thebibliography}

\end{document}
